\documentclass[11pt]{article}
\usepackage[utf8]{inputenc}
\usepackage{amsmath, amssymb}
\usepackage{geometry}
\usepackage{listings}
\usepackage{xcolor}
\usepackage{graphicx}
\usepackage{tabularx}
\usepackage{booktabs}
\usepackage{hyperref}

\geometry{a4paper, margin=1in}

\definecolor{codegray}{rgb}{0.5,0.5,0.5}
\definecolor{codepurple}{rgb}{0.58,0,0.82}
\definecolor{backcolour}{rgb}{0.95,0.95,0.92}

\lstset{
    language=Python,
    backgroundcolor=\color{backcolour},
    commentstyle=\color{codegray},
    keywordstyle=\color{magenta},
    stringstyle=\color{codepurple},
    basicstyle=\ttfamily\small,
    breaklines=true,
    showstringspaces=false
}

\title{Physics Lab: Quantum Wave Packets \& Scattering \\ \large Problem Set \& Implementation Guide}
\author{VB}
\date{\today}

\begin{document}

\maketitle

\section*{Problem 0: Start the engine}

Welcome to the Wave Packet simulator. Natural units ($\hbar=1, m=1$) are used everywhere to make our Fourier transforms clean.

\begin{center}
\renewcommand{\arraystretch}{1.5}
\begin{tabularx}{\textwidth}{@{} >{\bfseries\raggedright\arraybackslash}p{5.5cm} X @{}}
\toprule
\textcolor{blue}{Filename} & \textcolor{blue}{Purpose} \\ \midrule
\multicolumn{2}{@{}l}{\textbf{Direct Interaction}} \\ \midrule
\lstinline|implementation_tasks.py| & \textcolor{red}{Here you write your code.} Contains the core quantum mechanics solvers (wavefunctions, FFTs, and split-operator steps). You can add your own \lstinline|unittest.TestCase| classes at the bottom and run \lstinline|python implementation_tasks.py| to execute them. You can add your own \lstinline|unittest.TestCase| classes at the bottom and run \lstinline|python implementation_tasks.py| to execute them. \\
\lstinline|levels/level_*.py| & Interactive graphical sandbox for each specific problem. \textbf{Can be run standalone} with \lstinline|python levels/level_*.py| for visual debugging. \\
\lstinline|wave_explorer.py| & \textbf{The grand finale.} Run this when all levels are solved to launch the interactive tunneling demo. (Modes: \lstinline|barrier|, \lstinline|well|, \lstinline|finite_well|). \\ \midrule
\bottomrule
\end{tabularx}
\end{center}

\vspace{0.5cm}
Open \lstinline|implementation_tasks.py|. You will find empty un-implemented functions marked with \lstinline|pass|. Proceed sequentially through the problems below. Run the corresponding \lstinline|level_*.py| to debug your math. 


\section*{Problem 1: The Initial Packet}

\subsubsection*{Mathematical part}
\textbf{Given:} A position spatial grid $x$, a centroid $x_0$, a width $\sigma$, and a momentum $k_0$. \\
\textbf{Derive:} The normalized Gaussian wave packet. The formula is:
$$
\psi(x) = \left(\frac{1}{2\pi \sigma^2}\right)^{1/4} \exp\left(-\frac{(x-x_0)^2}{4\sigma^2}\right) e^{i k_0 x}
$$
Ensure you perform a discrete normalization pass: divide by $\sqrt{\sum |\psi|^2 \Delta x}$ to combat grid effects.

\subsubsection*{Implementation part}
\textbf{Implement:} \lstinline|gaussian_packet(x, x0, sigma, k0)|. Must return a complex \lstinline|numpy.array|. \\
\textbf{Test:} \lstinline|python levels/level_1_gaussian.py|. Does it look like a smooth bump? Does the integral display 1.000?


\section*{Problem 2: Fourier Duality}

\subsubsection*{Mathematical part}
\textbf{Given:} The spatial wavefunction $\psi(x)$ and grid spacing $\Delta x$. \\
\textbf{Derive:} The momentum-space wavefunction $\phi(k)$ and the conjugate $k$-grid. Use Fast Fourier Transform (FFT). Remember that $\phi(k) \approx \text{FFT}(\psi) \cdot \Delta x$. You must normalize $\phi(k)$ discretely.

\subsubsection*{Implementation part}
\textbf{Implement:} \lstinline|momentum_wavefunction(psi, dx)|. \\
\textbf{Test:} \lstinline|python levels/level_2_momentum_space.py|. Watch what happens when you adjust $\sigma$. As the packet gets narrow in space, does it broaden in momentum?


\section*{Problem 3: Free Evolution}

\subsubsection*{Mathematical part}
\textbf{Given:} The Hamiltonian for a free particle is purely kinetic: $\hat{H} = \hat{p}^2 / 2m$. In $k$-space, the propagator is an exact algebraic multiplication:
$$ \phi(k, t+\Delta t) = \phi(k, t) \exp\left(-i \frac{k^2 \Delta t}{2}\right) $$

\subsubsection*{Implementation part}
\textbf{Implement:} \lstinline|evolve_free_particle(psi, k, dt)|. Use FFT to go to $k$-space, apply the phase, and IFFT back. \\
\textbf{Test:} \lstinline|python levels/level_3_free_particle.py|. The wave packet should move to the right and spread.


\section*{Problem 4: Why Quantization?}

\subsubsection*{Mathematical part}
\textbf{Given:} An infinite square well from $x=0$ to $x=L$. \\
\textbf{Derive:} The eigenfunction $\psi_n(x) = \sqrt{\frac{2}{L}}\sin\left(\frac{n\pi x}{L}\right)$.

\subsubsection*{Implementation part}
\textbf{Implement:} \lstinline|well_eigenfunction(x, n, L)|. It returns the purely real function array. \\
\textbf{Test:} \lstinline|python levels/level_4_eigenfunctions.py|. \textbf{Critial Experiment:} move the $n$ slider to a non-integer. Notice how $\psi(L) \neq 0$. This explicitly shows why only integers are allowed — particles cannot exist inside infinite walls.


\section*{Problem 5: Split-Operator}

\subsubsection*{Mathematical part}
\textbf{Given:} The full potential $V(x)$. The Trotter formula splits the propagator:
$$ e^{-i \hat{H} \Delta t } \approx e^{-i \hat{V} \Delta t/2 } \, e^{-i \hat{T} \Delta t } \, e^{-i \hat{V} \Delta t/2 } $$

\subsubsection*{Implementation part}
\textbf{Implement:} \lstinline|split_operator_step(psi, k, V, dt)|. Apply the first half potential kick, hop to $k$-space, apply the full kinetic kick, hop to $x$-space, and apply the final half potential kick. (Use your code from Problem 3). \\
\textbf{Test:} \lstinline|python levels/level_5_split_operator.py|. Watch the packet scatter against the barrier... and then wrap around the edges!


\section*{Problems 6 \& 7: The Final Fixes}

You have two distinct problems to solve numerical artifacts:
\begin{itemize}
    \item \textbf{Level 6 (Infinite Well via DST):} Use the Discrete Sine Transform (DST) to exactly solve the infinite well boundary condition. \textbf{Implement:} \lstinline|dst_energy_levels(N, L)|.
    \item \textbf{Level 7 (Absorbing BCs):} In an open space, we want to kill scattering that reaches the grid boundary. \textbf{Implement:} \lstinline|absorbing_mask(N, gobble_frac)| using a $\sin^2$ taper envelope that goes to 0 near edges.
\end{itemize}


\section*{Problem 8: The Tunneling Explorer}

Once everything is complete, run:
\begin{lstlisting}[language=bash]
python wave_explorer.py --mode barrier
python wave_explorer.py --mode finite_well
\end{lstlisting}

\textbf{Task:} Experiment with the tunneling modes. Find the barrier height $V_0$ that cuts transmission down to exactly 1\%. Try the finite well and see resonant bound states reflecting back and forth.

\section*{\textcolor{red}{Submission}}
Submit \texttt{implementation\_tasks.py} and two screenshots proving your \lstinline|wave_explorer.py| achieved scattering and bound states.

\end{document}
